% GENERATED by scripts/build.py convert ch02 -- do not hand-edit; edit
% chapters/md/02_literature_review.md and regenerate.
% md sha256 581096fa2d7ea72f59f84fb70ef00b104aa691a21f6cab649858761c9bf14750
% canonical manuscript/v13_2026-10-17/02_literature_review.md sha256 581096fa2d7ea72f59f84fb70ef00b104aa691a21f6cab649858761c9bf14750
% md identical to canonical: True
% filter sha256 d487b0e18e456f08097065d7ceaaf39fef0daee29958941b5539c501e0e4c24a

\section{Literature review and conceptual framework}\label{sec:2}\ThesisCheckNumber{section}{2}

\subsection{From biological plausibility to demonstrated association}\label{sec:2.1}\ThesisCheckNumber{subsection}{2.1}

The hypothesised pathway from WASH to child growth runs through several distinct mechanisms, not one undisputed mechanism: inadequate water and sanitation can increase exposure to faecal--oral pathogens, raising diarrhoeal disease, soil-transmitted helminth infection, and sub-clinical environmental enteric dysfunction (EED), each potentially redirecting energy from linear growth, though EED specifically is the least settled \autocite{CummingCairncross2016}. Even the strongest recent evidence against basic WASH interventions reducing stunting does not challenge this plausibility: the authors of the three major trials discussed below state explicitly that ``the biological plausibility of WASH as public health interventions is not challenged by these findings'' \autocite{Cumming2019}, arguing the \emph{specific, basic} interventions tested were insufficient to translate the mechanism into a detectable effect. Black et al. \autocite*{Black2013} place WASH-related infection among the proximate determinants of undernutrition in their \emph{Lancet} framework, and Victora et al. \autocite*{Victora2010} document that growth faltering concentrates from birth to roughly 24 months. This review addresses a narrower question: whether a \emph{facility-type classification} --- an ``improved'' water source or sanitation facility in the WHO/UNICEF JMP sense used throughout --- shows a detectable statistical association with child growth once obvious confounders are accounted for, distinct from whether WASH-related infection can plausibly affect growth at all.

\subsection{Observational evidence and the confounding problem}\label{sec:2.2}\ThesisCheckNumber{subsection}{2.2}

A substantial observational literature finds WASH variables associated with child anthropometric outcomes in DHS and related data. Rakotomanana et al. \autocite*{Rakotomanana2020} link JMP water and sanitation indicators to linear growth in children 6--23 months across East African DHS surveys, in the expected direction. Addae et al. \autocite*{Addae2024}, using Ghana's 2017/18 MICS, find household sanitation and socioeconomic status jointly associated with stunting and wasting, and unprotected water associated with wasting but not stunting. Günther and Fink \autocite*{Gunther2010} pool 172 DHS surveys and find improved water and sanitation associated with lower child mortality; Spears \autocite*{Spears2013} finds cross-country height variation correlates with open-defecation rates.

The central difficulty this thesis faces directly is that WASH status is not randomly assigned: it correlates with wealth, education, and local infrastructure, which independently predict nutrition through channels unrelated to pathogen exposure. Headey and Palloni \autocite*{HeadeyPalloni2019} address this with a panel of 442 subnational regions across 59 countries, using within-region difference-in-differences. Their results are mixed: water access is insignificant for most outcomes (piped water alone predicts reduced stunting), and sanitation predicts large reductions in diarrhoea and mortality but \textbf{not} stunting or wasting once within-region design is used. This null sanitation--stunting finding, once region-level time-invariant confounding is netted out, parallels this thesis's own null household- and community-sanitation findings (Section 5); it is not evidence that WASH associations generally weaken under better controls, which would overstate what one study establishes.

\subsection{Credible causal evidence: what the major trials found}\label{sec:2.3}\ThesisCheckNumber{subsection}{2.3}

Randomised evidence comes primarily from three cluster-randomised trials, described individually because their exposure-timing structure differs materially from the cross-sectional literature above: WASH-Benefits Bangladesh, WASH-Benefits Kenya, and the SHINE trial in Zimbabwe enrolled pregnant women and followed children prospectively for 18--24 months from near birth --- unlike the single-survey-date measurement this thesis's own baseline design uses (Section 1.2). The three did not share one factorial design: Bangladesh and Kenya tested water, sanitation, hygiene, combined-WASH and nutrition as \textbf{separate arms}; SHINE used a simpler 2×2 design with a single combined WASH package. On linear growth, \textbf{every individual WASH-related arm in Bangladesh and Kenya, and SHINE's combined WASH arm, showed no statistically detectable effect} \autocite{Cumming2019}. The nutrition arm (no WASH) showed a separate, modest positive effect; combined WASH-plus-nutrition showed \textbf{no additional benefit over nutrition alone} in any trial. On diarrhoea, results were mixed: a substantial reduction in Bangladesh, none detectable in Kenya or Zimbabwe.

``No statistically detectable effect'' is not ``proven to be exactly zero'': these are adequately powered trials, but a null result reflects the trial's power, and the authors do not generalise the finding to every possible WASH intervention. The consensus statement attributes the null to the limited scope of the interventions tested --- household-level ``basic'' services, not ``safely managed'' networked services --- calling for research on more comprehensive WASH rather than concluding WASH cannot matter \autocite{Cumming2019}. This bears directly on this thesis: its facility-status exposure is likewise a classification, not a comprehensive intervention package, so a null or weak association here is consistent with, though it does not confirm, the trial finding, and equally consistent with confounding masking a real effect or this thesis's own limitations (Section 6).

A related quasi-experimental literature targets infant mortality. Geruso and Spears \autocite*{GerusoSpears2015} show Muslim infants in India survive to their first birthday at a higher rate than Hindu infants despite lower wealth, isolating a sanitation-specific channel from the fact that Hindus practise open defecation at a rate roughly 25 points higher nationally --- a religious difference unexplained by wealth, and a source of sanitation variation distinct from general local development that this thesis's own four-country comparison cannot replicate. The study is cited for this identification strategy, not as direct evidence about growth, which it does not examine.

\subsection{Household versus community exposure}\label{sec:2.4}\ThesisCheckNumber{subsection}{2.4}

Most DHS-based studies measure WASH at the household level only; a smaller strand examines community-level exposure. Donohue et al. \autocite*{Donohue2023} construct geospatially modelled community sanitation coverage and relate it to stunting in Nigeria and Zambia: a mixed result --- significant in Nigeria (adjusted OR = 0.97, 95\% CI {[}0.95, 0.99{]} per 10-point increase) after household sanitation lost significance, not significant in Zambia for either. Skoufias and Vinha \autocite*{SkoufiasVinha2026}, across twenty Sub-Saharan countries, find no significant water externality but a significant community-sanitation association with fever/diarrhoea and height-for-age --- contrasting with this thesis's own fully adjusted community HAZ models (Section 5.2), from a different sample and model. This thesis's leave-one-out construction follows the same logic; no study was found pairing it with a household-level comparison in the same four-country sample. Balk et al. \autocite*{Balk2004} provide an earlier precedent for DHS-cluster geography as a unit of analysis distinct from the household, for mortality rather than anthropometric outcomes.

A specific identification issue follows from this construction: the leave-one-out rate is a cluster-level property, so estimating it with cluster fixed effects would remove nearly all its substantive variation, identifying the coefficient almost entirely from the mechanical residual that arises because a different household is excluded from each child's own calculation (Section 4 derives this precisely). This thesis therefore uses region, not cluster, fixed effects for the community specification.

\subsection{Exposure timing and geospatial matching precedent}\label{sec:2.5}\ThesisCheckNumber{subsection}{2.5}

The single-survey-date limitation (Section 1.2) applies to this thesis's baseline design and to Rakotomanana et al.~and Addae et al.~above, but not to Günther and Fink (mortality outcomes) or Spears (a cross-country correlation), nor to the trials (Section 2.3) or Headey and Palloni's panel design. This is why the historical extension attempts a different, early-life-window exposure assignment within its own sample --- not a reason to expect comparing the two samples to isolate an exposure-timing effect (Section 6.2).

Assigning a gridded exposure surface to a survey respondent's location requires handling two documented DHS features: coordinates are randomly displaced for confidentiality \autocite{Burgert2013,PerezHeydrich2013}, so any extraction method reads pixels around the \emph{recorded}, not true, location; and a containing-pixel extraction (the single pixel a point falls in) differs from a buffer-based average. This thesis uses an area-weighted buffer average primarily and containing-pixel extraction as a labelled sensitivity (Section 3.6; Section 5.8). Johnson, Jacob, and Brown \autocite*{JohnsonJacobBrown2013} provide an earlier precedent linking satellite-derived surfaces to DHS clusters for a different outcome. Blom, Ortiz-Bobea, and Hoddinott \autocite*{Blom2022} match geocoded heat-exposure data to each child's developmental window across five West African countries --- the design structure this thesis's historical extension borrows, though their outcome is heat exposure's effect on nutrition, not WASH, and the study is cited only for its matching-method structure.

\subsection{This thesis's contribution and the gap it addresses}\label{sec:2.6}\ThesisCheckNumber{subsection}{2.6}

This thesis combines two elements not found together elsewhere in this review: household and leave-one-out community WASH compared directly within one sample and controls (Section 2.4); and a separate, bounded, preliminary attempt to assign modelled local WASH coverage to each child's early-life window via geospatial matching (Section 2.5), within the three historical-sample countries. The four-country combination alone is not evidence of substantive novelty; the claimed contribution is narrower --- a direct, consistently controlled comparison of the two exposure levels, and a transparent early-life exposure-matching attempt, including its own unresolved issues, that neither Donohue et al.~nor Skoufias and Vinha implement, since both examine only current-round community exposure. The gap that remains open, and that this thesis does not claim to close, is the one Headey and Palloni and the trial literature already identify at a larger scale: credibly separating a WASH-specific effect from the broader package of local development that improved facilities tend to accompany. The historical extension narrows only the exposure-timing dimension of that problem; it resolves neither the confounding dimension nor exposure-timing for the baseline sample, since the two samples differ in several other respects at once (Section 6.2).
